\subsection{Parity conditions, fair value and determination}

These entries cover what determines exchange rates: the parity conditions, fair value and forecasts, and payment flows, policy and crises. The market and its quote convention are in \hyperref[def:fx-market-real-exchange-rates]{the FX market}; cross and forward rates, bid--offer spreads, arbitrage and the value of a forward in \hyperref[meth:cross-forward-rates]{cross and forward rates} and the entries after it; regimes and capital flows in \hyperref[def:exchange-rate-regimes-capital-flows]{exchange-rate regimes}; and the carry trade in \hyperref[meth:carry-trade-uip]{carry and UIP}. Numbers are checked in \texttt{trading\_notes.numerical.fx\_mechanics\_parity.fx\_mechanics\_checks}.

\begin{qfresult}[International parity conditions: covered and uncovered interest parity, forward rates as forecasts, absolute, relative and ex ante purchasing power parity, the Fisher effect, real interest rate parity and the international Fisher effect, and how they connect]\label{res:international-parity-conditions}
\qffield{Statement} \citep[chap.~10, sec.~3.1, pp.~549--565]{economics-investment-decision-makers-2013}
\begin{itemize}
\item \emph{Covered interest parity}: the forward premium equals the interest differential, enforced by arbitrage (\hyperref[meth:cross-forward-rates]{CIP}).
\item \emph{Uncovered interest parity}: the expected change in the spot rate equals the interest differential, $\%\Delta S^e_{f/d}\approx i_f-i_d$; it requires risk neutrality and has no arbitrage to enforce it.
\item \emph{Forward rate parity}: if both CIP and UIP hold, the forward is an unbiased forecast of the future spot; asking whether it is unbiased ``is the same as asking whether uncovered interest rate parity holds'' (p.~553).
\item \emph{Purchasing power parity}: absolute (one price level across countries, which does not hold); relative (exchange-rate changes offset inflation differentials, $\%\Delta S_{f/d}\approx\pi_f-\pi_d$); ex ante (the same with expected inflation).
\item \emph{Fisher effect} and \emph{real interest rate parity}: nominal rate $\approx$ real rate $+$ expected inflation; if UIP and ex ante PPP hold, real rates converge across countries.
\item \emph{International Fisher effect}: the nominal interest differential equals the expected inflation differential, which follows from real interest rate parity and the Fisher effect (p.~561).
\end{itemize}
If all hold, the expected percentage change in the spot rate equals the forward premium, the nominal yield spread and the expected inflation spread, and the real yield spread is zero (p.~565).

\emph{Evidence}. CIP holds closely (with the post-2008 basis); UIP fails over short and medium horizons, so forwards are poor forecasts and high-yield currencies have not depreciated as predicted (pp.~553--556, 569). Relative PPP fails over one year but holds better over six years and more; ``Estimates place the half-life of PPP deviations at around three to five years'' (p.~560).

\qffield{Notation} $\pi$ inflation; $r$ real rate; f/d convention for the exchange rate.

\qffield{Conditions} The approximate forms drop cross products; exact forms use ratios of $(1+i)$ and $(1+\pi)$. The conditions relate expected changes and forwards to differentials; none pins down the level of the spot rate.

\qffield{Practice} \emph{Original.} With a 1.5 percent real rate everywhere and inflation of 3 percent at home and 1 percent abroad, the Fisher equation gives nominal rates of 4.545 and 2.515 percent; relative PPP and UIP then both predict the foreign currency appreciating 1.980 percent (d/f), equal to the forward premium under CIP.

\qffield{Pitfalls} Treating UIP as an arbitrage condition. Using PPP to time currencies over months. Confusing the forward discount of a high-yield currency with a forecast of depreciation.

\qffield{See also} \hyperref[meth:exchange-rate-forecasting-fair-value]{forecasting and fair value}; \hyperref[meth:carry-trade-uip]{carry trade}; \hyperref[ev:cip-deviations]{CIP deviations}.

\end{qfresult}

\begin{qfmethod}[Forecasting exchange rates and long-run fair value: forward, PPP, UIP and random-walk forecasts, real-interest-differential and risk-premium models, the IMF's equilibrium approaches, and shorter-term tools]\label{meth:exchange-rate-forecasting-fair-value}
\qffield{Statement} \emph{Evaluating forecasts} \citep[chap.~10, secs.~3.1--3.3 and~9, pp.~553--569 and~605--610]{economics-investment-decision-makers-2013}:
\begin{itemize}
\item the \emph{forward rate} and \emph{UIP} forecasts build in the interest differential, but because UIP fails over short horizons they are poor predictors; their errors are large relative to the predicted change;
\item \emph{relative PPP} is a weak guide over one year but an anchor over longer horizons, with deviations decaying slowly (half-life three to five years);
\item the \emph{random walk} (no change) is a hard benchmark to beat at short horizons;
\item fundamentals-based models explain long-run movements better than short-run ones, so shorter-term forecasting relies on technical analysis, order flow, sentiment and positioning (pp.~605--610).
\end{itemize}

\emph{Long-run fair value} (pp.~565--569). PPP is a fair-value framework over long horizons, but large, persistent departures make it a poor real-time gauge. The IMF's Consultative Group on Exchange Rate Issues used three approaches: the \emph{macroeconomic balance} approach (how far the exchange rate must move to close the gap between a country's medium-term current account and its sustainable level), the \emph{external sustainability} approach (the move needed to stabilise net foreign assets or debt at a chosen level), and a reduced-form econometric equilibrium real exchange rate. A simple model links the real rate to its long-run level, the real interest differential and relative risk premiums:
\[
q_{f/d}=\bar q_{f/d}+\big[(r_d-r_f)-(\phi_d-\phi_f)\big]\qquad\text{(eq.~10-12, p.~569)},
\]
so a currency's real value rises with an upward revision to its equilibrium, a higher relative real rate or a lower relative risk premium.

\qffield{Notation} $q$ real exchange rate (f/d), $\bar q$ its long-run equilibrium; $r$ real rates; $\phi$ risk premiums.

\qffield{Conditions} Equation (10-12) is EIDM's simplified form, whose two sides have different units unless the bracket is multiplied by $q$ (footnote~15, p.~569); fair-value estimates depend on the chosen price index, base year and model.

\qffield{Practice} \emph{Original.}
\begin{itemize}
\item \emph{Fair value from the model.} With $\bar q=1.00$, a real-rate gap of 1.5 points and a risk-premium gap of 0.8 points, $q=1.007$: the domestic currency sits 0.7 percent above its long-run value, supported by its real-rate advantage.
\item \emph{PPP fair value.} A base-year rate of 1.10 (d/f) with domestic prices up 10 percent and foreign up 5 percent gives 1.1524; a spot of 1.20 leaves the foreign currency 4.1 percent above PPP, a gap that would halve over three to five years on EIDM's evidence.
\item \emph{Forward against random walk.} On 600 simulated months in which a 3-point interest differential is not offset by depreciation, the random walk forecasts slightly better than the forward (root mean squared error 0.02543 against 0.02560): the forward's bias is small against monthly volatility of 2.5 percent, which is why tests need long samples.
\end{itemize}

\qffield{Pitfalls} Treating a fair-value gap as a short-term signal. Choosing the PPP base year that gives the desired answer. Judging forecasting skill on a few observations.

\qffield{See also} \hyperref[res:international-parity-conditions]{parity conditions}; \hyperref[meth:fx-valuation-signals]{FX valuation signals}; \hyperref[meth:fx-volatility-forecasts]{exchange-rate and volatility forecasts}.

\end{qfmethod}

\begin{qfdefinition}[Payment flows, policy and exchange rates: current and capital accounts, the Mundell--Fleming and monetary models, the policy mix, intervention and capital controls, and the warning signs of currency crises]\label{def:flows-policy-exchange-rates}
\qffield{Statement} \emph{Balance of payments} \citep[chap.~10, sec.~5, pp.~573--584]{economics-investment-decision-makers-2013}. The current account equals national saving less investment, private and public: $CA=(S_p-I)+(T-G)$. Persistent deficits tend to weaken a currency over the long run, through the need to attract capital and through the accumulation of external debt, while capital flows, being larger and faster, dominate in the short run; surges of inflows into emerging markets often end in sharp reversals.

\emph{Policy} (sec.~6, pp.~585--596). In the \emph{Mundell--Fleming} model with flexible rates and mobile capital, expansionary monetary policy lowers rates and depreciates the currency; expansionary fiscal policy raises rates and appreciates it if capital is mobile, but depreciates it if capital is immobile and the trade balance worsens. With fixed rates the central bank must offset monetary expansion (which becomes ineffective) and accommodate fiscal expansion. With high mobility, expansionary fiscal plus tight monetary policy is ``extremely bullish'' for a currency and the reverse bearish; with low mobility, policies act through trade flows, so uniformly tight policy is bullish (pp.~586--587, Exhibits~10-7 and~10-8). \emph{Monetary models} relate the rate to money supplies and expected inflation (with overshooting when prices are sticky); the \emph{portfolio balance} view holds that sustained deficits raise the supply of a country's bonds, requiring a weaker currency or higher yields unless policy reverses.

\emph{Intervention and controls} (sec.~7, pp.~597--601). Sterilised intervention leaves the money base unchanged; most studies find its effect in developed markets ``statistically insignificant or quantitatively unimportant'' because reserves are small relative to daily turnover, while emerging-market central banks, whose reserves are large relative to turnover in their currencies, are better placed (pp.~600--601). Capital controls are more effective against temporary than persistent flows and can preserve monetary independence while managing the exchange rate.

\emph{Crises} (sec.~8, pp.~602--605). Interest differentials, consensus forecasts and ratings failed to warn before the ERM, Mexican and Asian crises. An IMF study of ten variables around currency crises in 50 countries over 1975--1997, which EIDM summarises, found that before a crisis the real exchange rate is well above its tranquil-period mean, reserves fall sharply, the terms of trade deteriorate somewhat, inflation is higher, the ratio of M2 to reserves and broad money growth rise (money growth peaking about 18 months before), private credit grows quickly and equity prices boom and bust; the trade balance and real activity show no distinct pattern beforehand (pp.~603--604). One school holds that crises follow deteriorating fundamentals; another that sentiment shifts or contagion can strike sound economies, which limits any early-warning system (p.~603).

\qffield{Notation} $S_p$ private saving, $I$ investment, $T-G$ government balance, as shares of GDP.

\qffield{Conditions} The policy results assume sticky prices and given expectations; the crisis list is a set of stylised facts from studies EIDM summarises, and early-warning systems issue false signals.

\qffield{Practice} \emph{Original.}
\begin{itemize}
\item Private saving of 18 percent of GDP, investment of 22 percent and a budget deficit of 3 percent imply a current-account deficit of 7 percent, which must be financed by capital inflows of the same size.
\item A currency whose reserves fall 30 percent in five months while its real exchange rate rises 23 percent shows two of EIDM's warning signs; EIDM cautions that markets rarely price such crises in advance.
\end{itemize}

\qffield{Pitfalls} Reading a current-account deficit as a short-term sell signal. Expecting the same policy to move the currency the same way under high and low capital mobility. Relying on spreads or ratings for crisis warnings.

\qffield{See also} \hyperref[def:exchange-rate-regimes-capital-flows]{regimes and capital flows}; \hyperref[meth:monetary-fiscal-policy-mix]{policy mix}; \hyperref[def:currency-regimes-intervention]{intervention}; \hyperref[mod:margin-liquidity-spiral]{liquidity spirals}.

\end{qfdefinition}
